% === Begin Label Data ===
% ID: 686
% 难度星级: 
% 标签: 
% 备注: 
% 组卷引用次数: 0
% === End  Label Data ===

\begin{problem}{2008}{G}{江西卷（理）}{10}{立体几何}
连结球面上两点的线段称为球的弦. 半径为 $ 4 $ 的球的两条弦 $ AB $、$ CD $ 的长度分别等于 $ 2\sqrt{7} $、$ 4\sqrt{3} $, $ M $、$ N $ 分别为 $ AB $、$ CD $ 的中点，每条弦的两端都在球面上运动. 下列四个命题：
\begin{enumerate}
\item 弦 $ AB $、$ CD $ 可能相交于点 $ M $;
\item 弦 $ AB $、$ CD $ 可能相交于点 $ N $;
\item $ MN $ 的最大值为 $ 5 $;
\item $ MN $ 的最小值为 $ 1 $.
\end{enumerate}
其中真命题的个数为
\begin{choices}
\choice{{1 个}}
\choice{{2 个}}
\choice{{3 个}}
\choice{{4 个}}
\end{choices}
\end{problem}

\begin{answer}

\end{answer}

\begin{solutions}

\end{solutions}